 \section*{第 \ref{ch:b1:logic} 章——逻辑、集合与映射}

\begin{solution}{exo:b1:logic:1}
取否定，将 $\lnot$ 依次穿过各个量词
（\cref{prop:b1:logic:negquant}），并利用 $\lnot(P \implies Q) \iff
P \land \lnot Q$：
\begin{enumerate}
  \item $\exists x \in \R,\ \forall y \in \R,\ x + y \leq 0$;
  \item $\forall x \in \R,\ \exists y \in \R,\ xy \neq 0$;
  \item $\exists \varepsilon > 0,\ \forall \delta > 0,\ \exists x \in
        \R,\ \abs{x} \leq \delta \text{ 且 } \abs{f(x)} >
        \varepsilon$。
\end{enumerate}
命题（1）为真：对给定的 $x$，取 $y = -x + 1$；则 $x + y = 1 >
0$。命题（2）为真：$x = 0$ 满足 $xy = 0$，对所有 $y$ 成立。
\end{solution}

\begin{solution}{exo:b1:logic:2}
列真值表，以 T/F 表示四种情形 $(P, Q)$：
\begin{center}
\begin{tabular}{cc|c|c|c|c}
$P$ & $Q$ & $P \implies Q$ & $\lnot(P \implies Q)$ & $\lnot Q$ &
$P \land \lnot Q$ \\
\hline
T & T & T & F & F & F \\
T & F & F & T & T & T \\
F & T & T & F & F & F \\
F & F & T & F & T & F \\
\end{tabular}
\end{center}
第 $4$ 列与第 $6$ 列相同，由此证明等价。因此
“若函数可导，则它连续”的否定为：
“存在可导但不连续的函数”
（事实上这是假命题：原蕴含式为真，
见\cref{ch:b1:derivative}）。
\end{solution}

\begin{solution}{exo:b1:logic:3}
\emph{逆否证明。}设 $x < 1$。则 $x^3 < 1$（立方函数
单调递增），且 $x < 1$，故 $x^3 + x < 2$。这证明了
逆否命题，从而原命题成立。

\emph{反证法。}假设 $a > 0$ 是最小的严格正
实数。则 $a/2$ 为严格正数，且 $a/2 < a$（因 $a > 0$），
与最小性矛盾。故不存在这样的 $a$。
\end{solution}

\begin{solution}{exo:b1:logic:4}
\begin{enumerate}
  \item 归纳基础 $n = 0$：$2^0 = 1 = 2^1 - 1$。归纳步：设
        恒等式对 $n$ 成立，则
        \[
        \sum_{k=0}^{n+1} 2^k = (2^{n+1} - 1) + 2^{n+1}
        = 2 \cdot 2^{n+1} - 1 = 2^{n+2} - 1 .
        \]
  \item 归纳基础 $n = 0$：$4^0 + 5 = 6 = 3 \times 2$。归纳步：若
        $4^n + 5 = 3m$，则
        \[
        4^{n+1} + 5 = 4(4^n + 5) - 15 = 3(4m - 5),
        \]
        可被 $3$ 整除。
\end{enumerate}
\end{solution}

\begin{solution}{exo:b1:logic:5}
归纳步骤悄悄假设两组（“前
$n$ 支”与“后 $n$ 支”）有交集，使公共铅笔
把颜色从一组传到另一组。对 $n + 1 = 2$ 时两组
分别为 $\{$第一支铅笔$\}$ 与 $\{$第二支铅笔$\}$：两者不交，
论证随即失效。因此 $P(1) \implies P(2)$ 从未被证明，
归纳法随之失效——尽管 $P(n) \implies P(n+1)$ 对
每个 $n \geq 2$ 都成立。
\end{solution}

\begin{solution}{exo:b1:logic:6}
\begin{enumerate}
  \item $x \in A \setminus B \iff x \in A \land x \notin B \iff x \in
        A \land x \in \overline{B} \iff x \in A \cap \overline{B}$。
  \item 由（1）及分配律
        （\cref{prop:b1:logic:setalgebra}）：
        $(A \cup B) \cap \overline{C} = (A \cap \overline{C}) \cup
        (B \cap \overline{C})$。
  \item 设 $A \subseteq B$。则 $A \cup B \subseteq B$（两部分
        均在 $B$ 中），且恒有 $B \subseteq A \cup B$，故
        $A \cup B = B$。设 $A \cup B = B$，则 $A \cap B
        \subseteq A$ 恒成立，又 $A \subseteq A \cup B = B$ 给出
        $A \subseteq A \cap B$，故 $A \cap B = A$。设 $A \cap B =
        A$，则 $A = A \cap B \subseteq B$。三个条件
        由此等价（我们证明了一个蕴含循环）。
\end{enumerate}
\end{solution}

\begin{solution}{exo:b1:logic:7}
\begin{enumerate}
  \item $f$ 是\omterm{def:b1:logic:inj}{单射}（$n + 1 = m + 1 \implies n = m$），但不是
        \omterm{def:b1:logic:inj}{满射}：$0$ 在 $\N$ 中没有\omterm{def:b1:logic:map}{原像}。
  \item $g$ 是\omterm{def:b1:logic:inj}{双射}：$n \mapsto n - 1$ 是
        $\Z$ 上的双侧逆。
  \item $h$ 是\omterm{def:b1:logic:inj}{单射}：由 $\frac{x+1}{x-1} = \frac{x'+1}{x'-1}$ 得
        $(x+1)(x'-1) = (x'+1)(x-1)$，即 $xx' - x + x' - 1 = xx' -
        x' + x - 1$，故 $2x' = 2x$。它不是 $\R$ 上的\omterm{def:b1:logic:inj}{满射}：
        解 $y = \frac{x+1}{x-1}$ 得 $x(y - 1) = y + 1$，此方程
        当 $y = 1$ 时无解（方程化为 $0 = 2$）。
        若取值域为 $\R \setminus \{1\}$，同样的计算给出
        唯一的\omterm{def:b1:logic:map}{原像} $x = \frac{y+1}{y-1}$，故
        $h \colon \R \setminus \{1\} \to \R \setminus \{1\}$ 是
        \omterm{def:b1:logic:inj}{双射}，且 $h^{-1}(y) = \frac{y+1}{y-1} = h(y)$：$h$ 是
        自身的逆\omterm{def:b1:logic:map}{映射}。
\end{enumerate}
\end{solution}

\begin{solution}{exo:b1:logic:8}
\begin{enumerate}
  \item $x \in f^{-1}(B \cap B') \iff f(x) \in B \cap B' \iff f(x) \in
        B \land f(x) \in B' \iff x \in f^{-1}(B) \cap f^{-1}(B')$.
        关于像，$y \in f(A \cup A')$ 当且仅当 $y = f(x)$，其中 $x$
        属于 $A$ 或 $A'$，即 $y \in f(A)$ 或 $y \in f(A')$。
  \item 若 $y \in f(A \cap A')$，则 $y = f(x)$，其中 $x \in A$ 且
        $x \in A'$，于是 $y \in f(A)$ 且 $y \in f(A')$。下面说明这一包含可以是严格的：
        取 $f \colon \R \to \R$，$x \mapsto x^2$，$A = \{-1\}$，
        $A' = \{1\}$，则 $f(A \cap A') = f(\emptyset) = \emptyset$，
        而 $f(A) \cap f(A') = \{1\}$。
  \item （$\Leftarrow$）取 $A = \{x\}$，$A' = \{x'\}$，其中 $x \neq
        x'$，若 $f(x) = f(x')$，则 $f(A) \cap f(A') = \{f(x)\}$，
        而 $f(A \cap A') = \emptyset$，与所设
        等式矛盾，故 $f$ 为\omterm{def:b1:logic:inj}{单射}。（$\Rightarrow$）设 $f$ 为
        \omterm{def:b1:logic:inj}{单射}，且 $y \in f(A) \cap f(A')$，则 $y = f(x) = f(x')$，其中
        $x \in A$，$x' \in A'$，由\omterm{def:b1:logic:inj}{单射}性得 $x = x' \in A \cap
        A'$，故 $y \in f(A \cap A')$。结合（2）即得等号成立。
\end{enumerate}
\end{solution}

\begin{solution}{exo:b1:logic:9}
$g \circ f = \mathrm{id}_E$ 是\omterm{def:b1:logic:inj}{单射}且\omterm{def:b1:logic:inj}{满射}，故由
\cref{prop:b1:logic:comp}~(2)可知，$f$ 为\omterm{def:b1:logic:inj}{单射}，$g$ 为\omterm{def:b1:logic:inj}{满射}。
例：取 $E = \N$，$F = \Z$，$f$ 为包含\omterm{def:b1:logic:map}{映射} $n \mapsto n$，且
$g \colon \Z \to \N$，$g(n) = n$（$n \geq 0$），$g(n) = 0$（
$n < 0$）。于是 $g(f(n)) = n$ 对所有 $n \in \N$ 成立，但 $f$ 不是
\omterm{def:b1:logic:inj}{满射}，而 $g$ 不是\omterm{def:b1:logic:inj}{单射}。
\end{solution}

\begin{solution}{exo:b1:logic:10}
$x^2 - y^2 = x - y \iff (x - y)(x + y) = x - y \iff (x - y)(x + y - 1)
= 0 \iff y = x$ 或 $y = 1 - x$。\emph{自反性：}$y = x$ 成立。
\emph{对称性：}条件“$y = x$ 或 $y = 1 - x$”
        关于 $x$ 与 $y$ 对称（若 $y = 1 - x$，则 $x = 1 - y$）。
\emph{传递性：}设 $x \mathbin{\mathcal{R}} y$ 且 $y
\mathbin{\mathcal{R}} z$，逐一考察四种情形，每次 $z$ 都等于 $x$
或 $1 - x$（例如 $y = 1 - x$ 与 $z = 1 - y$ 给出 $z = x$）。
故 $\mathcal{R}$ 是\omterm{def:b1:logic:equiv}{等价关系}，且 $\mathrm{cl}(x) =
\{x,\, 1 - x\}$。该类恰有一个元素，当且仅当 $x = 1 - x$，
即 $x = \frac12$。
\end{solution}

\begin{solution}{exo:b1:logic:11}
设 $f \colon E \to \mathcal{P}(E)$ 为任一\omterm{def:b1:logic:map}{映射}，并令
$D = \{x \in E : x \notin f(x)\} \in \mathcal{P}(E)$。假设 $D =
f(a)$，其中 $a \in E$。若 $a \in D$，则由 $D$ 的定义，
$a \notin f(a) = D$，矛盾。若 $a \notin D$，则 $a \notin
f(a)$，由 $D$ 的定义得 $a \in D$，矛盾。因此 $D$
不在 $f$ 的像中，从而 $f$ 不是\omterm{def:b1:logic:inj}{满射}。（特别地，任何\omterm{def:b1:logic:sets}{集合}
都不与其\omterm{def:b1:logic:sets}{幂集}\omterm{def:b1:logic:inj}{双射}：
$\N$ 的子集个数比整数“更多”。）
\end{solution}

\begin{solution}{exo:b1:logic:12}
\begin{enumerate}
  \item （$\Rightarrow$）设 $f$ 为\omterm{def:b1:logic:inj}{满射}，且 $\Phi(B) =
        \Phi(B')$。对 $y \in B$，取 $x$ 使 $f(x) = y$，则 $x
        \in f^{-1}(B) = f^{-1}(B')$，于是 $y = f(x) \in B'$。因此 $B
        \subseteq B'$，同理 $B' \subseteq B$，故 $\Phi$ 为
        \omterm{def:b1:logic:inj}{单射}。（$\Leftarrow$）若 $f$ 不是\omterm{def:b1:logic:inj}{满射}，取 $y_0
        \in F$ 不在像中，则 $f^{-1}(\{y_0\}) = \emptyset =
        f^{-1}(\emptyset)$，且 $\{y_0\} \neq \emptyset$，故 $\Phi$ 不是
        \omterm{def:b1:logic:inj}{单射}。
  \item （$\Rightarrow$）设 $f$ 为\omterm{def:b1:logic:inj}{单射}，$A \subseteq E$。令
        $B = f(A)$，则 $f^{-1}(B) = \{x : f(x) \in f(A)\}$，且
        由\omterm{def:b1:logic:inj}{单射}性得 $f(x) \in f(A) \iff x \in A$，故 $\Phi(B) =
        A$，$\Phi$ 为\omterm{def:b1:logic:inj}{满射}。（$\Leftarrow$）若 $f$ 不是
        \omterm{def:b1:logic:inj}{单射}，取 $x \neq x'$ 使 $f(x) = f(x')$。每个
        \omterm{def:b1:logic:map}{原像}集 $f^{-1}(B)$ 含 $x$ 当且仅当它
        含 $x'$；故 $\{x\}$ 不是 $\Phi(B)$ 的形式，从而
        $\Phi$ 不是\omterm{def:b1:logic:inj}{满射}。
\end{enumerate}
\end{solution}

\begin{solution}{pb:b1:logic:1}
\textbf{1.} \emph{自反性：}$\mathrm{id}_E$ 是 $E$
到自身的\omterm{def:b1:logic:inj}{双射}。\emph{对称性：}若 $f \colon E \to F$ 为\omterm{def:b1:logic:inj}{双射}，
由\cref{thm:b1:logic:inverse}得 $f^{-1} \colon F \to E$，它也
是\omterm{def:b1:logic:inj}{双射}。\emph{传递性：}若 $f \colon E \to F$ 与 $g \colon F
\to G$ 均为\omterm{def:b1:logic:inj}{双射}，由\cref{prop:b1:logic:comp}~(1)可知 $g \circ f
\colon E \to G$ 为\omterm{def:b1:logic:inj}{双射}。（这只是“类似”\omterm{def:b1:logic:equiv}{等价关系}：
所有\omterm{def:b1:logic:sets}{集合}的全体本身不是\omterm{def:b1:logic:sets}{集合}，这一点由
\cref{exo:b1:logic:11}暗示的悖论可知；
真正重要的是这三条性质。）

\textbf{2.} 若 $f \colon E \to F$ 与 $g \colon F \to G$ 均为
\omterm{def:b1:logic:inj}{单射}，由\cref{prop:b1:logic:comp}~(1)可知 $g \circ f$ 为\omterm{def:b1:logic:inj}{单射}，故
$E \preceq G$。对第二点，把 $f$ 的值域缩小到其像，得
\omterm{def:b1:logic:map}{映射} $\tilde f \colon E \to f(E)$，$x \mapsto f(x)$，它由
$f(E)$ 的构造为\omterm{def:b1:logic:inj}{满射}，又由 $f$ 为\omterm{def:b1:logic:inj}{单射}，因而为\omterm{def:b1:logic:inj}{双射}：
$E \approx f(E)$。

\textbf{3.} （$\Rightarrow$）设 $f \colon E \to F$ 为\omterm{def:b1:logic:inj}{单射}，并
取定 $a \in E$（$E \neq \emptyset$）。定义 $s \colon F \to E$ 为：
$s(y)$ 为唯一的 $x$，使 $f(x) = y$，当 $y \in f(E)$ 时
（唯一性由\omterm{def:b1:logic:inj}{单射}性得），其余情形 $s(y) = a$。对每个 $x
\in E$，$s(f(x)) = x$，因而每个 $x$ 都被取到，$s$ 为\omterm{def:b1:logic:inj}{满射}。
（$\Leftarrow$）设 $s \colon F \to E$ 为\omterm{def:b1:logic:inj}{满射}。对每个 $x \in
E$，任取一个 $y_x \in F$ 使 $s(y_x) = x$，并令 $u(x) = y_x$。若
$u(x) = u(x')$，则 $x = s(u(x)) = s(u(x')) = x'$，故 $u \colon E \to F$
为\omterm{def:b1:logic:inj}{单射}。

\textbf{4.} $n \mapsto n + 1$ 把 $\N$ 映入 $\N^*$，且是\omterm{def:b1:logic:inj}{单射}
（$n + 1 = m + 1 \implies n = m$），且为\omterm{def:b1:logic:inj}{满射}（每个 $m \geq 1$ 都
可写成 $(m - 1) + 1$，其中 $m - 1 \in \N$）。对 $\sigma$，它把偶数
$0, 2, 4, \dots$ 映为 $0, 1, 2, \dots$，而奇数 $1, 3,
5, \dots$ 映为 $-1, -2, -3, \dots$。\omterm{def:b1:logic:inj}{单射}性：偶数输入的像落入
$\N$（$\sigma(n) = n/2 \geq 0$），奇数输入则落入
严格负整数（$\sigma(n) = -(n+1)/2 \leq -1$），因此两输入的像若相同，
则必在同一奇偶类内，而在其上 $\sigma$ 为
严格单调（$n/2 = m/2$ 或 $(n+1)/2 = (m+1)/2$ 推出 $n =
m$）。\omterm{def:b1:logic:inj}{满射}性：$k \geq 0$ 时可由 $\sigma(2k)$ 取到，$k \leq -1$ 时由
$\sigma(-2k - 1)$ 取到，其中 $-2k - 1 \geq 1$ 为奇数。于是 $\N \approx \N^*$，且
$\N \approx \Z$。

\textbf{5.} 由 $C_0 = E \setminus g(F) \subseteq C$，知若 $x \notin C$
则 $x \notin C_0$，即 $x \in g(F)$：存在 $y \in F$ 使
$g(y) = x$。若又有 $g(y') = x$，则由 $g$ 的\omterm{def:b1:logic:inj}{单射}性得 $y' = y$。
因此 $h$ 定义中的第二款确定了唯一的
良定元素 $g^{-1}(x)$。

\textbf{6.} 像集与并集可交换
（由\cref{exo:b1:logic:8}~(1)，先对 $f$ 再对 $g$ 应用）：
\[
g\bigl(f(C)\bigr)
= g\Bigl(f\Bigl(\bigcup_{n \in \N} C_n\Bigr)\Bigr)
= \bigcup_{n \in \N} g\bigl(f(C_n)\bigr)
= \bigcup_{n \in \N} C_{n+1}
= \bigcup_{n \geq 1} C_n \subseteq C .
\]

\textbf{7.} 设 $x \neq x'$ 属于 $E$。若二者都属于 $C$，则 $h(x) =
f(x) \neq f(x') = h(x')$，由 $f$ 的\omterm{def:b1:logic:inj}{单射}性。若两者都不在
$C$ 中，则 $g(h(x)) = x \neq x' = g(h(x'))$，故 $h(x) \neq h(x')$。若
$x \in C$ 而 $x' \notin C$（混合情形，互换记号后即可）：
设 $h(x) = h(x')$，即 $f(x) = g^{-1}(x')$。作用 $g$ 得
$x' = g(f(x)) \in g(f(C))$，由第 6 题得 $x' \in C$——
矛盾。于是各情形下均有 $h(x) \neq h(x')$，即 $h$ 为\omterm{def:b1:logic:inj}{单射}。

\textbf{8.} 设 $y \in F$。\emph{情形 1：$g(y) \notin C$。}则
$h(g(y)) = g^{-1}(g(y)) = y$，即元素 $g(y)$ 是一个\omterm{def:b1:logic:map}{原像}。
\emph{情形 2：$g(y) \in C$}，不妨设 $g(y) \in C_n$。由于 $g(y) \in
g(F)$，有 $g(y) \notin C_0 = E \setminus g(F)$，故 $n \geq 1$
且 $g(y) \in C_n = g(f(C_{n-1}))$，故存在 $x \in C_{n-1}$ 使
$g(y) = g(f(x))$。由 $g$ 的\omterm{def:b1:logic:inj}{单射}性得 $y = f(x)$，且 $x \in
C_{n-1} \subseteq C$，故 $h(x) = f(x) = y$。两种情形下 $y$ 均
被取到：$h$ 为\omterm{def:b1:logic:inj}{满射}，从而为\omterm{def:b1:logic:inj}{双射}。

\textbf{9.} 若 $E \preceq F$ 且 $F \preceq E$，取\omterm{def:b1:logic:inj}{单射} $f
\colon E \to F$ 与 $g \colon F \to E$；第 5--8 题构造了一个
\omterm{def:b1:logic:inj}{双射} $h \colon E \to F$，故 $E \approx F$。本命题并非
平凡，因为所给的两个\omterm{def:b1:logic:inj}{单射}彼此无关——二者
都不必是\omterm{def:b1:logic:inj}{满射}，且没有任何简单混合 $f$ 与 $g$ 的公式可直接定义\omterm{def:b1:logic:map}{映射}：
整个论证的关键在于把 $E$ \omterm{thm:b1:logic:partition}{划分}为区域 $C$
（在其上照搬 $f$）与其补集（在其上反向运用
$g$）。

\textbf{10.} (a) 包含\omterm{def:b1:logic:map}{映射} $\intoo01 \to \intcc01$ 为\omterm{def:b1:logic:inj}{单射}；
且 $x \mapsto \frac{x + 1}3$ 将 $\intcc01$ \omterm{def:b1:logic:inj}{单射}地映入
$\intcc{\frac13}{\frac23} \subseteq \intoo01$（它是仿射的，斜率
非零）。由第 9 题，$\intcc01 \approx \intoo01$——
\omterm{def:b1:logic:inj}{双射}，写出其显式表达式颇为繁琐。
(b) \emph{\omterm{def:b1:logic:inj}{单射}性。}设 $2^p(2q + 1) = 2^{p'}(2q' + 1)$，其中
不妨设 $p \leq p'$。两边除以 $2^p$，得 $2q + 1 = 2^{p' - p}(2q' + 1)$。
若 $p' > p$，则右边为偶数而左边为奇数——
不可能；故 $p = p'$，于是 $2q + 1 = 2q' + 1$，从而 $q = q'$。
\emph{\omterm{def:b1:logic:inj}{满射}性。}我们用强归纳法证明每个整数
$m \geq 1$ 都可写成 $2^p(2q + 1)$ 的形式。当 $m = 1$ 时，$p = q = 0$。
设 $m \geq 1$，并假设结论对 $\intint1m$ 中的
全体整数成立。若 $m + 1$ 为奇数，则 $m + 1 = 2q + 1$，取 $p = 0$。若
$m + 1$ 为偶数，则 $m + 1 = 2m'$，其中 $1 \leq m' \leq m$；由
归纳假设 $m' = 2^p(2q + 1)$，故 $m + 1 = 2^{p+1}(2q + 1)$。于是
$\varphi(p, q) = 2^p(2q + 1) - 1$ 取遍每个 $n \in \N$，且
$\varphi$ 为\omterm{def:b1:logic:inj}{双射} $\N \times \N \to \N$。

\textbf{11.} 由于 $A$ 为无限集，$A \setminus \{\varphi(0), \dots,
\varphi(n - 1)\}$ 永不为空，而
$\N$ 的最小元性质（用于证明 \cref{thm:b1:logic:induction}）使
上述递归定义合理。\emph{严格递增：}
$\varphi(n + 1)$ 属于 $A \setminus \{\varphi(0), \dots,
\varphi(n)\} \subseteq A \setminus \{\varphi(0), \dots, \varphi(n -
1)\}$，其最小元为 $\varphi(n)$；故 $\varphi(n + 1) \geq
\varphi(n)$，又等号不成立，从而 $\varphi(n+1) >
\varphi(n)$。\emph{$\varphi(n) \geq n$：}由归纳法，$\varphi(0)
\geq 0$，且 $\varphi(n + 1) \geq \varphi(n) + 1 \geq n + 1$。
\emph{\omterm{def:b1:logic:inj}{单射}性}由严格单调性推出。
\emph{到 $A$ 上的\omterm{def:b1:logic:inj}{满射}性：}假设某个 $a \in A$ 从不
被取到。由于 $\varphi(a + 1) \geq a + 1 > a$，考虑所有 $n$，其中
$\varphi(n) > a$；这样的整数构成的\omterm{def:b1:logic:sets}{集合}非空，设 $n$ 为其最小元。对每个
$k < n$，有 $\varphi(k) \leq a$，因而 $\varphi(k) < a$（$a$ 未
被取到）。于是 $a$ 属于 $A \setminus \{\varphi(0), \dots,
\varphi(n - 1)\}$，且 $a < \varphi(n)$，与
$\varphi(n)$ 定义中的极小性相矛盾。于是 $\varphi$ 为 $\N \to A$ 的\omterm{def:b1:logic:inj}{双射}，从而
$A$ 可数。

\textbf{12.} 设 $E \preceq \N$，记\omterm{def:b1:logic:inj}{单射}为 $f$，则 $E \approx
f(E)$（第 2 题）。若 $f(E)$ 为有限集，则 $E$ 为有限集；若 $f(E)$
为无限集，由第 11 题得 $f(E) \approx \N$，故由
传递性（第 1 题）得 $E \approx \N$。反之，有限集与可数集显然可\omterm{def:b1:logic:inj}{单射}入 $\N$ 中。
捷径是：由 $E \preceq \N$ 及 $\N
\preceq E$ 即可直接
由康托–施罗德–伯恩斯坦定理得 $E \approx \N$——无需枚举论证。

\textbf{13.} 设 $f \colon E \to \N$ 与 $g \colon F \to \N$ 均为
\omterm{def:b1:logic:inj}{单射}。则 $(x, y) \mapsto \varphi\bigl(f(x), g(y)\bigr)$ 是一个
\omterm{def:b1:logic:inj}{单射} $E \times F \to \N$：若两像相同，由 $\varphi$ 的\omterm{def:b1:logic:inj}{单射}性（第 10 题）
得 $f(x) = f(x')$ 且 $g(y) = g(y')$，
由\omterm{def:b1:logic:inj}{单射}性即得 $x = x'$，$y = y'$。对于 $\Z \times \N^*$，两个因子均
可数（第 4 题），故 $\Z \times \N^* \preceq \N$；它
为无限集（它包含 $\{0\} \times \N^*$），故由
第 12 题，它可数。

\textbf{14.} 每个有理数 $r$ 有唯一表示 $r =
p/q$，其中 $p \in \Z$，$q \in \N^*$，且分数既约
（唯一性在\cref{ch:b1:arith}中证明；当 $r = 0$ 时取
$0/1$）。\omterm{def:b1:logic:map}{映射} $r \mapsto (p, q)$ 是\omterm{def:b1:logic:inj}{单射}：数对
确定 $r = p/q$。因此 $\Q \preceq \Z \times \N^* \preceq \N$
由第 13 题。由 $\N \subseteq \Q$ 得 $\N \preceq \Q$，
由第 12 题（或直接由康托–施罗德–伯恩斯坦定理）得 $\Q
\approx \N$：有理数集可数。

\textbf{15.} 对每个 $n$，取定\omterm{def:b1:logic:inj}{单射} $f_n \colon E_n \to \N$。
对 $x \in \bigcup_n E_n$，令 $n(x)$ 为\emph{最小}的 $n$，使得
$x \in E_n$，并令 $u(x) = \varphi\bigl(n(x), f_{n(x)}(x)\bigr)
\in \N$。若 $u(x) = u(x')$，由 $\varphi$ 的\omterm{def:b1:logic:inj}{单射}性得 $n(x) =
n(x') = n$，且 $f_n(x) = f_n(x')$，从而
$x = x'$，由 $f_n$ 的\omterm{def:b1:logic:inj}{单射}性。于是该并集\omterm{def:b1:logic:inj}{单射}入 $\N$，从而至多可数。

\textbf{16.} 设 $\Psi(F) = \sum_{i \in F} 2^i$，其中 $F \subseteq \N$
为有限集（$\Psi(\emptyset) = 0$）。假设 $F \neq F'$，令 $m$ 为两者相异的最大元素，
不妨设 $m \in F \setminus
F'$（必要时交换两者名字）。$> m$ 的元素要么同时属于两者，要么同时不属于，
故对两和式的贡献相等；比较
元素 $\leq m$ 的贡献：
\[
\sum_{i \in F,\, i \leq m} 2^i \geq 2^m
> 2^m - 1 = \sum_{k=0}^{m-1} 2^k
\geq \sum_{i \in F',\, i \leq m} 2^i ,
\]
利用\cref{exo:b1:logic:4}中的等比求和。由此 $\Psi(F)
\neq \Psi(F')$，$\Psi$ 为\omterm{def:b1:logic:inj}{单射}，而
$\N$ 的有限子集全体构成的\omterm{def:b1:logic:sets}{集合}至多可数；它是无限的（包含所有单点集），
从而可数。

\textbf{17.} 将 $A \subseteq \N$ 映到其示性函数 $\mathbf 1_A
\colon \N \to \{0,1\}$，$\mathbf 1_A(n) = 1$（若 $n \in A$），$0$
（否则）；将 $u \in \{0,1\}^{\N}$ 映到 $A_u = \{n \in \N : u(n) =
1\}$。两个\omterm{def:b1:logic:map}{映射}互为逆\omterm{def:b1:logic:map}{映射}：$A_{\mathbf 1_A} = A$ 且
$\mathbf 1_{A_u} = u$（在每个 $n$ 处验证即可）。由
\cref{thm:b1:logic:inverse}，两者都是\omterm{def:b1:logic:inj}{双射}，故 $\mathcal{P}(\N)
\approx \{0,1\}^{\N}$。

\textbf{18.} 对每个 $n$，$d(n) = 1 - \Phi(n)(n) \neq \Phi(n)(n)$，
因此数列 $d$ 与 $\Phi(n)$ 在指标 $n$ 处不同，即 $d \neq
\Phi(n)$。因此任何 $\Phi$ 都不是\omterm{def:b1:logic:inj}{满射}，由第 3 题，
也不存在\omterm{def:b1:logic:inj}{单射} $\{0,1\}^{\N} \to \N$，故 $\{0,1\}^{\N}$ 不是
至多可数。通过第 17 题中的对应关系，\omterm{def:b1:logic:map}{映射} $\Phi
\colon \N \to \{0,1\}^{\N}$ 即为\omterm{def:b1:logic:map}{映射} $f \colon \N \to
\mathcal{P}(\N)$，而 $d$ 对应\omterm{def:b1:logic:sets}{集合} $D = \{n : n
\notin f(n)\}$（事实上 $d(n) = 1 \iff \Phi(n)(n) = 0 \iff n \notin
f(n)$）：对角线论证\emph{正是}康托在
\cref{exo:b1:logic:11}中对 $E = \N$ 的证明。

\textbf{19.} 将 $x_n = 0.d_1(n)\,d_2(n)\,d_3(n)\dots$ 写成规范
形式，并定义 $\delta_n = 5$（若 $d_n(n) \neq 5$），$\delta_n = 6$
（若 $d_n(n) = 5$），则令 $x = 0.\delta_1\delta_2\delta_3\dots$，该
展开式只使用数字 $5$ 与 $6$，故不以
全为 $9$ 的尾部结束：它是某个实数 $x \in \intco01$ 的规范展开。
对每个 $n$，第 $n$ 位数字在 $x$ 与 $x_n$ 中不同
（由构造，$\delta_n \neq d_n(n)$），而规范展开
是唯一的，故 $x \neq x_n$。因此没有数列能穷尽 $\intco01$：由
第 3 题再次可知，$\intco01$ 不是至多可数的。

\textbf{20.} 由 $\intco01 \subseteq \R$，若存在\omterm{def:b1:logic:inj}{单射} $\R \to \N$，
则其限制给出 $\intco01$ 上的\omterm{def:b1:logic:inj}{单射}，与第 19 题矛盾：
$\R$ 不可数。若 $\R \setminus \Q$ 至多可数，
则 $\R = \Q \cup (\R \setminus \Q)$ 将是两个至多可数集的并，
从而由第 15 题，它至多可数（取
$E_0 = \Q$，$E_n = \R \setminus \Q$（$n \geq 1$））——
矛盾。因此无理数不可数。确切地说：
在 $\R$ 之中，有理数构成可数集，而它们的
补集不可数；任何\omterm{def:b1:logic:inj}{双射}都不能将 $\R
\setminus \Q$ 与 $\Q$ 一一对应——即无理数严格多于有理数，
尽管两者都无限且都稠密。

\textbf{21.} $p/q$（$q \neq 0$）是多项式 $qX - p$ 的根，它是一个
非零整系数多项式。$\sqrt 2$ 是
$X^2 - 2$ 的根。对 $x = \sqrt 2 + \sqrt 3$，$x^2 = 5 + 2\sqrt 6$，故
$x^2 - 5 = 2\sqrt 6$，且 $(x^2 - 5)^2 = 24$，即\
\[
x^4 - 10x^2 + 1 = 0 :
\]
$\sqrt 2 + \sqrt 3$ 是 $X^4 - 10X^2 + 1$ 的根。

\textbf{22.} 将 $P = a_0 + a_1X + \dots + a_nX^n$（次数 $\leq n$，
整系数）映到 $(a_0, \dots, a_n) \in \Z^{n+1}$：这是
\omterm{def:b1:logic:inj}{单射}，因为多项式由其系数唯一确定。由
对 $n$ 归纳：$\Z^1 = \Z$ 可数（第 4 题），且
$\Z^{n+2} \approx \Z^{n+1} \times \Z$ 至多可数，由
第 13 题。故次数有界的整系数多项式全体构成的\omterm{def:b1:logic:sets}{集合}
至多可数；它是无限的（包含所有常数），由第 12 题，
故可数。

\textbf{23.} 全体整系数多项式构成的\omterm{def:b1:logic:sets}{集合}为 $\bigcup_{n \in
\N} \{P : \deg P \leq n,\ P \text{ 具有整系数}\}$，它是
可数集的可数并：由
第 15 题知至多可数，又无限，故可数。

\textbf{24.} 对每个非零整系数多项式 $P$，根集
$R_P = \{x \in \R : P(x) = 0\}$ 是有限集（至多 $\deg P$
个元素，在本阶段不加证明地承认）。由第 23 题，非零整系数多项式
可以排成 $P_0, P_1, P_2, \dots$；于是 $\mathcal{A} =
\bigcup_{n \in \N} R_{P_n}$ 是有限（从而至多可数）集的可数并，
由第 15 题至多可数。它包含 $\Q$（第 21 题），
故为无限集，于是 $\mathcal{A}$
可数。

\textbf{25.} 若 $\R \setminus \mathcal{A}$ 至多可数，
则 $\R = \mathcal{A} \cup (\R \setminus \mathcal{A})$ 至多可数
（第 15 题），与第 20 题矛盾。因此
\omterm{pb:b1:logic:1}{超越数}存在，而且它们构成不可数集，
而\omterm{pb:b1:logic:1}{代数数}——其中包括由整数经根式构造的每一个数——
仅构成一个可数的“骨架”，
位于 $\R$ 之中。论证结构概要：第 1--3 题
建立了比较的语言；康托--施罗德--伯恩斯坦（第
5--9 题）以两个简单的\omterm{def:b1:logic:inj}{单射}代替一个巧妙的\omterm{def:b1:logic:inj}{双射}来证明等势，
并用于证明 $\intcc01 \approx \intoo01$，
用于 $\Q$，并贯穿第五部分；配对\omterm{def:b1:logic:inj}{双射}（第
10 题）支撑了乘积与可数并（第 13、15 题），
进而支撑了 $\Q$、整系数多项式与
$\mathcal{A}$ 的讨论；对角线论证（第 18--19 题）提供了
唯一的严格不等式 $\N \prec \R$，使整个故事
变得不平凡。康托的结论在哲学上极具震撼力：
证明中并未给出任何\omterm{pb:b1:logic:1}{超越数}，却表明在
等势的意义下，\emph{几乎所有}实数都是
\omterm{pb:b1:logic:1}{超越数}。具体指出一个\omterm{pb:b1:logic:1}{超越数}——$\pi$ 或
$\eu$——需要完全不同的数学，且又
多花了数十年的功夫。
\end{solution}
